% Fig. 2: diagrama de flujo tipo PRISMA 2020. Todos los números son macros de generado/conteos.tex.
\begin{tikzpicture}[node distance=2.6mm and 2.4mm]
  % Identificación
  \node[caja=3.75cm] (bases) {\textbf{Bases de datos} (2019--2026)\\
     arXiv: Q1 \nArxivQuno{} $+$ Q2 \nArxivQdos{} $=$ \nArxivTotal\\
     OpenAlex: Q1 \nOaQuno{} $+$ Q2 \nOaQdos{} $=$ \nOaTotal};
  \node[caja=2.45cm,right=of bases] (otros) {\textbf{Otras fuentes}\\ clásicos, citas y
     búsqueda dirigida: \nOtros};
  % Tamizaje
  \node[caja=3.75cm,below=4mm of bases] (tam) {\textbf{Tamizaje por título y resumen}\\
     \nTamizados{} registros (sin \nDup{} duplicados)};
  \node[excl,right=of tam] (extam) {Excluidos (E1--E5): \nExclTitulo};
  \node[caja=3.75cm,below=of tam] (tc) {\textbf{Evaluados en elegibilidad}: \nPasan};
  \node[excl,right=of tc] (extc) {Excluidos tras leer: \nExclTexto};
  \node[caja=3.75cm,below=of tc,fill=udazul!10] (inc) {\textbf{Incluidos: \nTotalIncluidos}
     (\nIncluidosBases{} de bases $+$ \nOtros{} de otras fuentes)\\ \nRecientes{} desde 2021;
     \nMatriz{} en la Tabla~\ref{tab:matriz}};
  \draw[flecha] (bases) -- (tam);
  \draw[flecha] (tam) -- (tc);
  \draw[flecha] (tc) -- (inc);
  \draw[flecha] (otros.east) -- ++(1.6mm,0) |- (inc.east);
  \draw[flecha,udplata] (tam.east) -- (extam.west);
  \draw[flecha,udplata] (tc.east) -- (extc.west);
  % Etiquetas de fase
  \foreach \n/\t in {bases/Ident., tam/Tamiz., inc/Incl.}
    \node[fase,rotate=90,anchor=south] at ([xshift=-1.2mm]\n.west) {\t};
\end{tikzpicture}
